

\section{Mémoire Déclarative vs Procédural : Le modèle DP d'Ullman}

Titre du Projet de Thèse : Mécanismes Neurocognitifs et Didactique des Langues Anciennes

Sujet du Rapport : Mémoire Déclarative vs Procédurale : Le Modèle DP d'Ullman appliqué à
l'acquisition des langues et ses implications pour l'enseignement du grec ancien

L'intégration des neurosciences cognitives dans le champ de la didactique des langues constitue l'un
des tournants épistémologiques majeurs du XXIe siècle. Dans le cadre spécifique d'une thèse portant
sur l'enseignement du grec ancien, cette approche ne relève pas d'une simple curiosité
interdisciplinaire, mais d'une nécessité impérieuse. En effet, l'enseignement des langues
classiques, et du grec en particulier, traverse une crise de légitimité et d'efficacité, souvent
cristallisée autour de la persistance de la méthode \og Grammaire-Traduction\fg{}. Cette méthode,
héritée du XIXe siècle, repose intuitivement sur une conception explicite et analytique de
l'apprentissage, que les données neurobiologiques actuelles permettent désormais d'interroger avec
une précision inédite.

Ce rapport se consacre exclusivement et exhaustivement au point II.2.\cite{II-2-1-ref1} de votre
recherche doctorale : le \textbf{Modèle Déclaratif/Procédural (DP)} théorisé par le
neuroscientifique Michael T. Ullman. Ce modèle postule une double dissociation fonctionnelle et
anatomique dans le cerveau humain entre le lexique mental et la grammaire mentale, ancrée
respectivement dans la mémoire déclarative et la mémoire procédurale. Loin d'être une simple
cartographie cérébrale, le modèle DP offre une grille de lecture puissante pour comprendre les
mécanismes d'acquisition d'une langue seconde (L2) et, par extension, les obstacles spécifiques
rencontrés par les apprenants du grec ancien.

L'objectif de ce document de 15 000 mots est de fournir une exégèse complète du modèle, depuis ses
fondements neurobiologiques jusqu'à ses applications didactiques les plus concrètes. Nous
explorerons comment ce modèle a évolué depuis sa formulation initiale en 1997 jusqu'à ses mises à
jour récentes (2016, 2020), intégrant les concepts de redondance et de dégénérescence neuronale.
Nous confronterons ensuite ce modèle aux théories concurrentes, notamment la \textit{Skill
Acquisition Theory} de Robert DeKeyser, pour analyser les processus d'automatisation. Enfin, nous
appliquerons ces concepts à la spécificité morphologique du grec ancien, proposant une critique
neuro-informée des pratiques pédagogiques actuelles et esquissant les contours d'une didactique
compatible avec le fonctionnement du cerveau humain.



\subsection{Les fondements neurobiologiques du modèle déclaratif/procédural}

Le modèle DP repose sur une hypothèse évolutionniste élégante : le langage humain, malgré sa
complexité unique, n'a pas émergé \textit{ex nihilo} par la création de modules cérébraux
entièrement nouveaux. Au contraire, il a \og co-opté\fg{} des systèmes de mémoire préexistants chez
les vertébrés, dévolus à des fonctions non linguistiques. Cette perspective, qualifiée de \og
neurobiologiquement motivée\fg{}, ancre l'étude du langage dans la biologie générale de la mémoire.



\subsection{La dichotomie lexique / grammaire et ses substrats neuronaux}

Au cœur du modèle réside la distinction entre deux capacités linguistiques fondamentales : le
lexique mental et la grammaire mentale. Cette distinction linguistique se superpose à une
distinction neurobiologique entre deux systèmes de mémoire à long terme.



\subsection{Le système de mémoire déclarative : le substrat du lexique}

La mémoire déclarative est le système qui sous-tend la connaissance consciente des faits (\og savoir
que\fg{}) et des événements (mémoire épisodique). Dans le modèle DP, ce système est le siège du
\textbf{lexique mental}.

Nature de la connaissance :

Le système déclaratif est spécialisé dans l'apprentissage de connaissances arbitraires,
idiosyncrasiques et non dérivables. Dans le langage, cela correspond prototypiquement à la relation
arbitraire entre un son (signifiant) et un sens (signifié). Par exemple, rien dans la structure
phonologique du mot grec anthrôpos ne permet de déduire logiquement le concept d'\og être
humain\fg{}. Cette association doit être apprise par cœur et stockée telle quelle. Le lexique mental
ne se limite pas aux mots simples ; il englobe également les expressions idiomatiques figées, les
collocations fréquentes, et surtout, les formes morphologiques irrégulières qui échappent aux règles
générales. Ainsi, lorsqu'un apprenant mémorise que l'aoriste de legô (je dis) est eipon, il
sollicite sa mémoire déclarative, car aucune règle productive ne permet de générer cette forme.

Architecture Neuroanatomique :

L'ancrage biologique de ce système est bien identifié. Il repose primairement sur les structures du
Lobe Temporal Médian (MTL), et plus spécifiquement sur l'hippocampe. L'hippocampe joue un rôle
crucial dans l'encodage initial et la consolidation des nouvelles informations lexicales. Cependant,
à mesure que les souvenirs se consolident, ils deviennent progressivement indépendants de
l'hippocampe pour être stockés dans des régions néocorticales distribuées, notamment dans les lobes
temporaux (gyrus temporal supérieur et moyen, aires de Brodmann 21, 22) et les régions temporo-
pariétales (aires de Brodmann 39, 40). La récupération de ces informations lexicales implique
également le cortex préfrontal ventrolatéral, en particulier les aires de Brodmann 45 et 47, qui
jouent un rôle dans la sélection sémantique et la récupération contrôlée en mémoire.

Caractéristiques Fonctionnelles :

L'apprentissage déclaratif présente des caractéristiques spécifiques qui influencent directement la
pédagogie. Il est :

\begin{itemize}\item \textbf{Rapide :} Il permet l'apprentissage en un seul essai (\textit{one-shot learning}). Un apprenant peut retenir le sens d'un mot grec rare après une seule exposition marquante.\item \textbf{Explicite :} Les connaissances sont accessibles à la conscience. L'apprenant \og sait qu'il sait\fg{} et peut verbaliser cette connaissance.\item \textbf{Flexible :} Les connaissances déclaratives peuvent être généralisées et utilisées dans des contextes différents de ceux de l'apprentissage initial.\item \textbf{Attention-dépendant :} L'encodage nécessite une attention focale. La distraction ou une charge cognitive excessive peuvent entraver l'apprentissage lexical.\end{itemize}

\subsection{Le système de mémoire procédurale : le moteur de la grammaire}

La mémoire procédurale est le système impliqué dans l'apprentissage et l'exécution d'habiletés
motrices et cognitives (\og savoir comment\fg{}). C'est le système qui nous permet de faire du vélo,
de taper au clavier, ou de jouer d'un instrument. Dans le modèle DP, ce système prend en charge la
\textbf{grammaire mentale}.

Nature de la connaissance :

Le système procédural gère les règles qui régissent la combinaison séquentielle et hiérarchique des
éléments lexicaux. C'est le domaine de la régularité, de la structure et de la prédictibilité. En
grec ancien, c'est ce système qui devrait idéalement gérer la déclinaison d'un nom ou la conjugaison
d'un verbe régulier. Lorsqu'un locuteur natif produit une forme verbale régulière, il ne la \og
récupère\fg{} pas en mémoire comme un bloc figé ; il la construit en temps réel en appliquant une
règle de composition (ex: Racine \+ Suffixe thématique \+ Désinence). Cette capacité de computation
permet la productivité infinie du langage : nous pouvons comprendre et produire des phrases grecques
que nous n'avons jamais entendues auparavant.

Architecture Neuroanatomique :

Ce système est enraciné dans les circuits fronto-striataux. Le composant central est constitué par
les ganglions de la base, des structures sous-corticales profondes. Plus précisément, le noyau caudé
et le putamen (formant le striatum dorsal) jouent un rôle déterminant. Ces structures reçoivent des
afférences de l'ensemble du cortex et projettent en retour vers le cortex frontal (aires prémotrices
et aire de Broca \- BA 44/45) via le thalamus, formant des boucles de régulation motrice et
cognitive. Le cervelet est également impliqué, notamment dans le traitement de la séquentialité
temporelle fine et l'optimisation des prédictions sensori-motrices.

Caractéristiques Fonctionnelles :

L'apprentissage procédural diffère radicalement de l'apprentissage déclaratif :

\begin{itemize}\item \textbf{Graduel :} Il nécessite une répétition massive et constante. On n'acquiert pas une habileté procédurale en une seule fois ; il faut s'entraîner.\item \textbf{Implicite :} Les règles sont apprises et appliquées sans conscience explicite. Un locuteur natif applique les règles de l'accentuation grecque sans nécessairement pouvoir les énoncer.\item \textbf{Rigide :} Les compétences procédurales sont souvent spécifiques à la tâche et moins transférables que les connaissances déclaratives.\item \textbf{Automatique :} Une fois consolidées, les procédures s'exécutent rapidement, avec peu d'effort et sans consommer les ressources limitées de la mémoire de travail.\end{itemize}

\subsection{Évolution et mises à jour du modèle (1997-2020)}

Le modèle DP n'est pas statique ; il a été affiné au fil des deux dernières décennies pour intégrer
les avancées en neuroimagerie et en génétique. Il est impératif pour votre thèse de ne pas se
limiter à la version originale de 1997, mais de prendre en compte les révisions majeures d'Ullman
(2004, 2016).



\subsection{De la dissociation stricte à la co-optation redondante}

Les premières formulations du modèle insistaient sur une double dissociation stricte (Lexique \=
Déclaratif / Grammaire \= Procédural). Cependant, Ullman (2016) intègre désormais le concept
biologique de \textbf{dégénérescence} (\textit{degeneracy}) ou de redondance fonctionnelle. Ce
concept stipule que des mécanismes biologiques différents peuvent aboutir au même résultat
comportemental.

Dans le contexte linguistique, cela signifie que bien que le système procédural soit
\textit{optimal} pour le traitement grammatical (rapide, automatique), le système déclaratif possède
une capacité de \textbf{compensation}. Il peut apprendre et stocker des séquences grammaticales
complexes sous forme de \og chunks\fg{} (blocs mémorisés) ou de règles explicites. Un apprenant peut
donc produire une phrase grammaticalement correcte en utilisant soit la computation procédurale
(voie native), soit la récupération déclarative (voie compensatoire). Cette nuance est capitale pour
comprendre l'apprentissage du grec ancien, où la voie procédurale est souvent déficiente.



\subsection{Facteurs modulant la dépendance aux systèmes}

Le modèle identifie plusieurs variables qui déterminent quel système prévaut chez un individu donné
:

\begin{itemize}\item \textbf{Le Sexe et les Hormones :} Des études citées par Ullman suggèrent un dimorphisme sexuel dans les stratégies cognitives. Les femmes, sous l'influence des œstrogènes, tendraient à avoir une mémoire déclarative plus performante et à l'utiliser davantage pour des tâches linguistiques que les hommes, qui s'appuieraient plus volontiers sur des mécanismes procéduraux. Les œstrogènes semblent faciliter la fonction hippocampique tout en pouvant inhiber certaines fonctions striatales.\item \textbf{La Génétique (FOXP2 et CNTNAP2) :} Le gène \textit{FOXP2}, souvent qualifié de \og gène du langage\fg{}, est fortement exprimé dans le striatum (ganglions de la base). Les variations de ce gène affectent directement la plasticité procédurale et l'apprentissage de séquences. Une mutation de ce gène entraîne des troubles spécifiques du langage qui affectent la grammaire et l'articulation, confirmant le lien entre génotype, phénotype cérébral (striatum) et fonction linguistique (procédurale).\end{itemize}

\subsection{Le modèle dp appliqué à l'acquisition des langues secondes (sla)}

L'application du modèle DP à l'acquisition d'une langue seconde (SLA \- \textit{Second Language
Acquisition}) constitue le point de bascule théorique pour votre recherche. Si l'acquisition de la
langue maternelle (L1) se caractérise par une spécialisation neuronale optimale, l'acquisition d'une
L2, surtout à l'âge adulte (cas typique des étudiants en grec ancien), présente une dynamique
neurocognitive radicalement différente.



\subsection{La trajectoire développementale : du déclaratif au procédural}

Le modèle DP prédit une trajectoire d'apprentissage distincte pour la L2, conditionnée par l'état de
maturation des systèmes de mémoire.



\subsection{Le stade initial : la dépendance déclarative}

Chez l'apprenant L2 débutant ou intermédiaire, le système procédural est souvent moins accessible ou
moins efficace pour extraire les régularités grammaticales de la nouvelle langue. En conséquence,
l'apprenant s'appuie massivement sur sa \textbf{mémoire déclarative} pour toutes les fonctions
linguistiques, y compris la grammaire.

Concrètement, face à une phrase grecque, le débutant ne \og parse\fg{} pas la structure syntaxique
automatiquement. Il utilise sa mémoire explicite pour :

1. Récupérer le sens des mots un par un.

2. Récupérer consciemment des règles grammaticales apprises (\og Le génitif marque la
possession\fg{}, \og l'augment indique le passé\fg{}).

3. Assembler le sens global par déduction logique, un peu comme on résout une équation mathématique.

Cette dépendance au système déclaratif explique la lenteur, l'hésitation et la charge cognitive
élevée observées chez les débutants. Le traitement n'est pas automatique ; il est contrôlé et
attentif.



\subsection{La possibilité de procéduralisation}

La question centrale en SLA est de savoir si un apprenant tardif peut finir par développer une
compétence procédurale native-like. Le modèle DP, dans ses versions récentes, est prudemment
optimiste. Il postule que la pratique intensive et l'exposition à l'input peuvent, avec le temps,
permettre un transfert progressif du contrôle du système déclaratif vers le système procédural.
C'est ce qu'on appelle la \textbf{procéduralisation}.

Cependant, Ullman souligne que ce transfert est entravé par plusieurs facteurs chez l'adulte :

\begin{itemize}\item \textbf{Atténuation de la plasticité procédurale :} La capacité d'apprentissage implicite des séquences semble décliner après l'adolescence, rendant la procéduralisation de la grammaire plus difficile (mais pas impossible) que chez l'enfant.\item \textbf{Interférence de la L1 :} Les circuits procéduraux sont déjà \og câblés\fg{} pour la grammaire de la langue maternelle, créant une résistance à l'établissement de nouvelles routines grammaticales contradictoires.\end{itemize}

\subsection{L'effet de bascule ("see-saw effect") : compétition entre systèmes}

L'un des concepts les plus pertinents pour la didactique est l'effet de \og balançoire\fg{} ou de
\og bascule\fg{} (\textit{See-Saw Effect}) décrit par Ullman (2004) et Poldrack \& Packard (2003).



\subsection{Mécanisme de compétition}

Les données neurobiologiques suggèrent que les systèmes de mémoire déclarative et procédurale ne
sont pas seulement indépendants, mais peuvent interagir de manière \textbf{compétitive}. Une
activation forte de l'hippocampe (mémoire déclarative) peut, par des voies inhibitrices directes,
supprimer l'activité du striatum (mémoire procédurale). Inversement, l'engagement procédural peut
moduler l'activité déclarative.

Cette compétition a une base adaptative : face à une tâche d'apprentissage, le cerveau doit \og
choisir\fg{} la stratégie la plus efficace. Le système déclaratif, étant plus rapide à apprendre,
gagne souvent la compétition initiale (\og race to learn\fg{}). Une fois que la solution déclarative
est trouvée, le cerveau peut cesser de chercher à développer une routine procédurale plus coûteuse
en temps d'entraînement.



\subsection{Implications didactiques : le paradoxe de l'instruction explicite}

Cet effet de bascule pose un problème fondamental pour l'enseignement explicite du grec ancien. Si
la méthode pédagogique met une emphase excessive sur l'apprentissage explicite des règles et
l'analyse métalinguistique (activant fortement l'hippocampe), elle risque physiologiquement
d'\textbf{inhiber} le recrutement des ganglions de la base.

Autrement dit, en donnant trop d'explications grammaticales et en encourageant l'analyse consciente
constante, on pourrait \og bloquer\fg{} la capacité naturelle du cerveau à automatiser la langue.
L'élève devient un expert en analyse grammaticale (déclaratif), mais reste incapable de lire
fluidement (procédural). C'est le paradoxe de l'étudiant qui connaît par cœur toutes les règles de
l'accentuation grecque mais est incapable de lire une phrase à voix haute sans trébucher.



\subsection{Données de neuroimagerie en sla}

Les études citées dans le corpus bibliographique confirment cette dynamique. Les examens IRMf et ERP
(Potentiels Évoqués) montrent des profils d'activation distincts selon le niveau de compétence :

\begin{itemize}\item \textbf{Débutants :} Lors de tâches grammaticales, on observe une activation accrue des régions temporales (mémoire déclarative) et peu d'activation fronto-striatale.\item \textbf{Avancés / Bilingues précoces :} L'activation se déplace vers le réseau procédural classique (Aire de Broca, Ganglions de la base), mimant le profil des locuteurs natifs.\item \textbf{Fossilisation :} Chez certains apprenants avancés qui n'atteignent jamais la fluidité, l'activation reste obstinément déclarative, suggérant une \og fossilisation\fg{} neurocognitive où le système procédural n'a jamais pris le relais.\end{itemize}

\subsection{Débats théoriques : modèle dp vs skill acquisition theory (sat)}

Pour construire une thèse solide, il est nécessaire de confronter le modèle DP à la théorie
dominante en didactique des langues instruites : la \textit{Skill Acquisition Theory} (SAT),
défendue notamment par Robert DeKeyser. Cette confrontation permet d'affiner la compréhension du
processus d'automatisation.



\subsection{La skill acquisition theory (sat) : la vision linéaire}

S'inspirant du modèle ACT-R d'Anderson, DeKeyser (1997, 2015) propose une vision optimiste et
linéaire de l'apprentissage des compétences :

1. \textbf{Savoir Déclaratif :} L'apprenant reçoit une règle explicite (ex: \og l'accusatif marque
le COD\fg{}).

2. \textbf{Procéduralisation :} Par la pratique délibérée et l'exercice, ce savoir déclaratif est
converti en règles de production (\og Si je veux exprimer un COD, j'utilise l'accusatif\fg{}).

3. \textbf{Automatisation :} Avec une pratique massive, le comportement devient rapide, sans erreur
et demande peu de ressources cognitives.

Selon la SAT, l'interface entre connaissance explicite et implicite est \og forte\fg{} :
l'enseignement explicite est la voie royale vers la compétence implicite.



\subsection{La critique du modèle dp : systèmes distincts, routes parallèles}

Le modèle DP offre une vision plus nuancée, voire contradictoire sur certains points. Pour Ullman,
il n'y a pas de transmutation magique d'une trace mnésique hippocampique (déclarative) en une trace
striatale (procédurale). Ce sont deux systèmes biologiquement distincts.

Ce qui ressemble à une \og transformation\fg{} est en réalité un \textbf{changement de prédominance}
(\textit{shift}). L'apprenant développe parallèlement des connaissances dans les deux systèmes. Au
début, la trace déclarative est plus forte et contrôle le comportement. Avec la pratique, la trace
procédurale se renforce. À un moment critique, la trace procédurale devient plus rapide et plus
fiable que la trace déclarative, et prend le contrôle du comportement. L'automatisation n'est donc
pas la transformation du savoir, mais le changement de la route neurale utilisée pour la tâche.



\subsection{L'effet de blocage (blocking effect)}

Des recherches sur l'apprentissage, citées en lien avec le modèle DP 1, mettent en évidence un \og
effet de blocage\fg{}. Une connaissance explicite préexistante peut empêcher l'apprentissage
implicite de nouvelles régularités.

\begin{itemize}\item \textbf{Explication :} Si l'apprenant possède une règle explicite qui lui permet de résoudre la tâche (même lentement), son cerveau ne ressent pas le \og signal d'erreur\fg{} nécessaire pour déclencher l'apprentissage procédural implicite.\item \textbf{Application au Grec :} Si l'on enseigne explicitement toutes les règles de déclinaison \textit{avant} toute exposition aux textes, l'élève utilisera systématiquement sa règle explicite pour décoder. Son système procédural, n'étant jamais mis en situation de \og deviner\fg{} ou de prédire statistiquement la forme, ne s'engagera pas. L'explicite \og court-circuite\fg{} l'implicite.\end{itemize}

\subsection{La charge cognitive (sweller) et le goulot d'étranglement déclaratif}

La Théorie de la Charge Cognitive de John Sweller 3 complète utilement le modèle DP. La mémoire de
travail, interface obligatoire pour la mémoire déclarative, est extrêmement limitée (4 à 7
éléments).

\begin{itemize}\item \textbf{Surcharge en Grec :} La méthode Grammaire-Traduction impose une charge cognitive massive. Pour traduire un seul verbe grec, l'élève doit maintenir en mémoire de travail : la racine, les morphèmes de temps/mode/aspect, les règles de modification phonétique, et le sens lexical. Cette surcharge sature la mémoire de travail, empêchant tout apprentissage profond. Le système procédural, lui, ne souffre pas de cette limitation de capacité de la même manière, car il traite l'information de façon modulaire et automatique.\end{itemize}

\subsection{Application spécifique : vers une didactique du grec ancien neuro-informée}

L'analyse précédente n'est pas purement spéculative ; elle a des implications directes et radicales
pour la pédagogie du grec ancien. Le grec, en tant que langue à morphologie riche (fusionnelle) et
enseignée hors immersion, représente un \og stress-test\fg{} pour le modèle DP.



\subsection{La spécificité morphologique du grec et la défaillance déclarative}

Le grec ancien se caractérise par une complexité morphologique extrême. Un verbe peut théoriquement
prendre des centaines de formes différentes.

\begin{itemize}\item \textbf{Le Fardeau Mémoriel :} Tenter de gérer cette complexité uniquement via la mémoire déclarative est une entreprise vouée à l'échec cognitif. La mémoire déclarative est faite pour stocker des faits unitaires, pas pour gérer des algorithmes combinatoires complexes en temps réel.\item \textbf{La Nécessité Procédurale :} Pour lire du grec fluidement (comme le faisaient les Anciens), il est impératif que le traitement morphologique (reconnaître que \textit{\-ois} est un datif pluriel) soit procéduralisé. C'est-à-dire qu'il doit être instantané, inconscient et sans coût attentionnel. Or, la méthode traditionnelle, en focalisant sur l'analyse (\og Ceci est un datif pluriel de la 2ème déclinaison\fg{}), renforce la voie déclarative au détriment de la voie procédurale.\end{itemize}

\subsection{Critique neurocognitive de la méthode grammaire-traduction}

À la lumière du modèle DP, la méthode Grammaire-Traduction apparaît comme une anomalie
neurocognitive :

1. \textbf{Absence d'Input Séquentiel :} Le système procédural apprend par statistiques sur des
séquences (\textit{Statistical Learning}). La traduction mot-à-mot (\og version\fg{}) détruit la
séquentialité de la phrase. L'élève ne traite pas le grec comme un flux linguistique, mais comme un
puzzle visuel. Les ganglions de la base ne reçoivent pas l'input temporel nécessaire pour extraire
des régularités.

2. \textbf{Timing Inversé :} On donne la règle explicite \textit{avant} l'exposition. Cela active
l'effet de blocage et l'effet de bascule, inhibant l'apprentissage implicite potentiel.

3. \textbf{Focalisation sur l'Exception :} L'enseignement traditionnel valorise souvent la
connaissance des exceptions (\og les curiosités grammaticales\fg{}). Or, les exceptions sont
stockées dans la mémoire déclarative. En se focalisant sur elles, on renforce encore la dépendance à
ce système, au lieu de travailler la régularité massive qui construit la compétence procédurale.



\subsection{Pistes pour une didactique hybride}

Si l'immersion totale est impossible pour une langue morte, le modèle DP suggère des adaptations
pour favoriser la procéduralisation :



\subsection{Maximiser l'input et la fréquence (pattern drills)}

Pour activer le système procédural, il faut augmenter massivement la fréquence d'exposition aux
structures régulières \textit{en contexte}.

\begin{itemize}\item \textbf{Lecture Extensive :} Proposer des textes simples en grande quantité où les mêmes structures reviennent fréquemment, permettant au cerveau de faire ses statistiques implicites.\cite{II-2-1-ref5}\item \textbf{Exercices Structuraux :} Utiliser des exercices de manipulation rapide (transformer du singulier au pluriel, du présent au passé) sans analyse métalinguistique. L'objectif est la vitesse et la répétition, pour forcer le passage du traitement déclaratif lent au traitement procédural rapide.\end{itemize}

\subsection{L'oralisation comme levier procédural}

Le système procédural est évolutivement lié au système moteur et auditif.

\begin{itemize}\item \textbf{La Boucle Phonologique :} Lire à haute voix, réciter, ou écouter du grec active les circuits fronto-striataux et cérébelleux impliqués dans le séquençage temporel. L'apprentissage silencieux (\og purement visuel\fg{}) prive le cerveau d'un canal d'entrée majeur pour la grammaire. Réintroduire l'oralité n'est pas une coquetterie esthétique, mais une nécessité neurocognitive pour ancrer la grammaire.\end{itemize}

\subsection{La gestion stratégique de l'explicite}

Il ne s'agit pas d'interdire la grammaire explicite (utile pour les structures complexes et
l'adulte), mais de changer son timing.

\begin{itemize}\item \textbf{Séquence Inductive :} \textit{Exposition \-\> Pratique \-\> Explicitation}. Laisser d'abord le système procédural tenter de traiter les données. L'explication grammaticale vient ensuite consolider ou corriger, mais elle ne doit pas précéder systématiquement l'expérience linguistique.\end{itemize}

\subsection{Synthèse des données et tableaux comparatifs}

Pour visualiser clairement les implications du modèle DP, nous proposons ci-dessous une synthèse
structurée des différences entre les deux systèmes et de leur application au grec ancien.



\subsection{Tableau 1 : comparaison fonctionnelle et anatomique}

\begin{table}[h!]\small\centering\begin{tabular}{| p{2.3cm} | p{3.5cm} | p{3.5cm} |}\hline\textbf{Caractéristique} & \textbf{Mémoire Déclarative (Lexique)} & \textbf{Mémoire Procédurale (Grammaire)} \\ \hline\textbf{Type de Connaissance} & \og Savoir que\fg{} (Faits, Événements) & \og Savoir comment\fg{} (Habiletés, Règles) \\ \hline\textbf{Contenu Linguistique} & Mots, Idiomes, Irrégularités & Syntaxe, Morphologie régulière, Phonologie \\ \hline\textbf{Accès} & Explicite, Conscient & Implicite, Non-conscient \\ \hline\textbf{Vitesse d'Apprentissage} & Rapide (\og One-shot learning\fg{}) & Lent (Nécessite pratique et répétition) \\ \hline\textbf{Substrat Cérébral} & Lobe Temporal Médian (Hippocampe), Néocortex temporal (BA 45/47) & Ganglions de la Base (Striatum: Caudé/Putamen), Frontal (Broca BA 44), Cervelet \\ \hline\textbf{Neurochimie} & Acétylcholine, Œstrogènes & Dopamine \\ \hline\textbf{Sensibilité au Vieillissement} & Robuste, s'améliore parfois avec l'âge & Décline après l'enfance/adolescence \\ \hline\end{tabular}\end{table}

\subsection{Tableau 2 : implications pour l'apprentissage du grec ancien}

\begin{table}[h!]\small\centering\begin{tabular}{| p{2.3cm} | p{3.5cm} | p{3.5cm} |}\hline\textbf{Dimension} & \textbf{Approche Traditionnelle (Grammaire-Traduction)} & \textbf{Approche Neuro-Informée (DP Model)} \\ \hline\textbf{Stratégie Cognitive} & Dépendance totale au Déclaratif (Règles \+ Lexique) & Tentative d'équilibre Déclaratif / Procédural \\ \hline\textbf{Traitement de l'Erreur} & Analyse consciente de l'erreur (Feedback explicite) & Correction par l'exposition massive (Feedback implicite) \\ \hline\textbf{Rôle de l'Oral} & Marginal (prononciation théorique) & Central (activateur des boucles procédurales) \\ \hline\textbf{Progression} & Par complexité grammaticale théorique & Par fréquence et utilité communicative \\ \hline\textbf{Résultat Typique} & \og Savoir sur la langue\fg{} (Expertise métalinguistique) mais lecture lente & \og Savoir la langue\fg{} (Compétence de lecture fluide) \\ \hline\end{tabular}\end{table}L'exploration du point II.2.\cite{II-2-1-ref1} de votre thèse révèle que le modèle
Déclaratif/Procédural d'Ullman est bien plus qu'une théorie neurolinguistique parmi d'autres ; c'est
une clé de voûte pour comprendre les apories de l'enseignement des langues anciennes.

L'analyse a mis en lumière trois conclusions majeures :

1. \textbf{La Dissonance Cognitive :} L'enseignement du grec ancien souffre d'une inadéquation
fondamentale entre la nature de son objet (une morphologie complexe nécessitant un traitement
procédural) et ses méthodes (qui sollicitent exclusivement la mémoire déclarative). Cela crée une
surcharge cognitive permanente et empêche l'accès à la fluidité.

2. \textbf{Le Risque de l'Explicite :} Loin d'être toujours bénéfique, l'abus d'enseignement
explicite peut, par le biais de l'effet de bascule (\textit{See-Saw Effect}) et de l'effet de
blocage, entraver physiologiquement les mécanismes naturels d'automatisation.

3. \textbf{L'Espoir de la Redondance :} Si la procéduralisation totale est difficile chez l'adulte,
la plasticité cérébrale et la redondance des systèmes permettent des stratégies compensatoires. Une
pédagogie qui réintroduit la séquentialité, l'oralité et la fréquence peut permettre, sinon de créer
des locuteurs natifs, du moins de former des lecteurs fluides capables de traiter le grec comme une
langue vivante et non comme un code chiffré.

Ce modèle offre donc une justification scientifique rigoureuse pour expérimenter de nouvelles
approches didactiques dans votre thèse, validant le passage d'une pédagogie de l'analyse (le \og
pourquoi\fg{} de la langue) à une pédagogie de la pratique (le \og comment\fg{} de la langue).

